% Notas de aula — Capítulo 13: Ciclones Extratropicais
% Curso: Meteorologia Prática (baseado em R. Stull, Practical Meteorology, CC BY-NC-SA 4.0)
% Caixas verdes = conteúdo literal do quadro; linhas verdes = encenação.
\documentclass[11pt]{article}
\usepackage{../comum/preambulo}
\graphicspath{{../figuras/cap13/}}

\title{Capítulo 13 --- Ciclones Extratropicais\\
{\large Notas de aula (roteiro de quadro-negro)}}
\author{Curso de Meteorologia Prática}
\date{}

\begin{document}
\maketitle
\thispagestyle{fancy}

\noindent\textbf{Plano sugerido:} 2 aulas de 2\,h.
\emph{Aula 1:} \S\ref{sec:vida}--\S\ref{sec:eady} (ciclo de vida
norueguês, vorticidade, instabilidade baroclínica).
\emph{Aula 2:} \S\ref{sec:spinup}--\S\ref{sec:sintese} (spin-up,
tendência de pressão, jet streaks, autodesenvolvimento, bombas).

\section{O protagonista das latitudes médias}\label{sec:vida}

\begin{noquadro}[abertura]
\textbf{Cap.~13 --- Ciclones Extratropicais}\\[2pt]
o ciclone é a \emph{máquina} que executa o serviço do Cap.~11:\\
transporta calor para o polo e momento para o jato,\\
usando a energia guardada no contraste térmico (baroclinia).
\end{noquadro}

Tudo o que o curso montou converge aqui: a frente que ondula
\javimos{12}{frentes}, a onda de Rossby que a embala
\javimos{11}{ondas de Rossby}, a convergência que levanta
\javimos{10}{continuidade} e a chuva frontal que aquece a coluna
\javimos{7}{precipitação}.

\quadro{Desenhar a sequência do modelo norueguês em 4 quadros: frente
estacionária $\to$ ondulação $\to$ ciclone maduro com setor quente
$\to$ oclusão. Anotar os tempos: $\sim$1--2 dias por quadro.}

O \textbf{ciclo de vida norueguês} (escola de Bergen
\javimos{12}{Bergen}): nasce como ondulação numa frente, amadurece com
setor quente bem definido, oclui quando a fria alcança a quente
\javimos{12}{oclusões} e morre em 5--8 dias. No HS tudo gira no sentido
\emph{horário} em torno da baixa — treinar os desenhos com o sul para
baixo.

\section{Vorticidade: a moeda do ciclone}\label{sec:vort}

\begin{formadif}[equação da vorticidade (T1)]
Tomando o rotacional das equações do movimento
\javimos{10}{equações do movimento}:
\[ \frac{D(\zeta + f)}{Dt} \simeq -(\zeta + f)\,
\nabla_h\!\cdot\vec U_h, \]
o termo de \textbf{esticamento}: convergência encolhe a coluna
horizontalmente e ela gira mais rápido (bailarina fechando os braços).
No HS, $\zeta$ do ciclone tem o mesmo sinal de $f$ (negativo… mas
convém trabalhar com $|\zeta|$ crescendo — fixar convenção no quadro).
\end{formadif}

\begin{noquadro}[onde nasce o giro]
convergência $D<0$ sobre coluna com $(\zeta + f)$\\
$\Rightarrow$ $|\zeta|$ cresce — o ciclone é filho da convergência\\
e a convergência de superfície é imposta pela \textbf{altitude}
(advecção de vorticidade, jet streak)
\end{noquadro}

O Caso~1 do notebook calcula $\zeta$ e sua advecção numa onda de
500\,hPa: a região a leste do cavado é a ``zona de sucção'' onde o
ciclone de superfície nasce e para onde ele se move — a regra de bolso
de qualquer carta de altitude \javimos{9}{cartas de 50\,kPa}.

\begin{formadif}[vorticidade potencial e a ciclogênese de sotavento]
Para uma coluna adiabática de espessura $\Delta z$, conserva-se a
\textbf{vorticidade potencial}
\[ \frac{\zeta + f}{\Delta z} = \text{const}. \]
É o esticamento em forma de lei de conservação
\javimos{11}{tipos de vorticidade}. Aplicação imediata: ar de oeste
cruzando uma cordilheira é \emph{achatado} no topo (perde $\zeta$,
curva anticiclônico) e \emph{esticado} na descida — ganha $\zeta$
ciclônica e forma o cavado de sotavento: a \textbf{ciclogênese de
sotavento} (Stull \S13.3) dos Alpes, das Rochosas e da Cordilheira
dos Andes (as baixas do Chaco e os ciclones que nascem a leste da
cordilheira e amadurecem no Prata).
\end{formadif}

\section{Por que existem ciclones? Instabilidade baroclínica}
\label{sec:eady}

\quadro{Desenhar as isentrópicas inclinadas da zona baroclínica e uma
troca de parcelas dentro do ``cone de instabilidade'': a fria desce e a
quente sobe — energia potencial liberada.}

\begin{formalivro}[taxa de Eady]
\[ \sigma_{max} \simeq 0{,}31\,\frac{|f|}{N}\,
\frac{\partial U}{\partial z},\qquad
\lambda_{max} \simeq 3{,}9\,\frac{N\,H}{|f|} \simeq 4000\ \mathrm{km}. \]
O cisalhamento vem do vento térmico \javimos{11}{vento térmico}; $N$ é
a estratificação \javimos{5}{Brunt--Väisälä}; $NH/|f|$ é o raio de
deformação de Rossby (T2).
\end{formalivro}

\begin{noquadro}[os números que explicam a carta]
$\lambda_{max} \simeq 4000$\,km $=$ o tamanho dos ciclones\\
e-dobra $\sim$1--2 dias no inverno $=$ a velocidade com que o tempo
muda\\
inverno (contraste maior) $=$ temporada de ciclones (N2)
\end{noquadro}

O ciclone é o modo instável da zona baroclínica: extrai energia
potencial do contraste equador--polo e a converte em vento — cumprindo
o transporte que a célula de Hadley não alcança
\javimos{11}{limite da Hadley}.

\section{Spin-up e morte: a EDO do ciclone}\label{sec:spinup}

\begin{formadif}[esticamento × Ekman (Caso~3)]
\[ \frac{d\zeta}{dt} = -(f+\zeta)\,D - \frac{\zeta}{\tau_E},
\qquad \tau_E \sim \text{2--4 dias}. \]
Com $D<0$ imposto pela altitude, o crescimento é
super-exponencial; o bombeamento de Ekman \javimos{10}{espiral de
Ekman} põe o teto e, quando a divergência de altitude desliga, mata o
ciclone em $\sim\tau_E$ (N3). Vida completa: 5--8 dias.
\end{formadif}

\subsection{Tendência de pressão}

\begin{formadif}[por que a pressão cai (T3)]
$P_{sfc}$ é o peso da coluna \javimos{1}{massa da coluna}:
\[ \frac{\partial P_{sfc}}{\partial t} =
-g\int_0^\infty \nabla\cdot(\rho\vec U)\,dz. \]
O ciclone aprofunda quando a coluna \emph{exporta} massa — divergência
em altitude vencendo a convergência de superfície. Aquecimento da
coluna (calor latente da chuva frontal \javimos{7}{calor latente})
também derruba $P_{sfc}$: umidade acelera bombas.
\end{formadif}

\section{Jet streaks e autodesenvolvimento}\label{sec:jet}

\quadro{Desenhar o streak com os 4 quadrantes e marcar (HS!) entrada
polar e saída equatorial como quadrantes de divergência em altitude.}

\begin{formalivro}[os quatro quadrantes (T4; Caso~4)]
Na entrada do streak o ar acelera; o vento ageostrófico transversal
$v_{ag} = (1/f)\,DU/Dt$ monta duas células de circulação: divergência
em altitude na \textbf{entrada polar} e na \textbf{saída equatorial}
(HS) — os berços preferenciais de ciclogênese.
\end{formalivro}

\textbf{Autodesenvolvimento}: o ciclone cava sua própria
intensificação — a advecção quente à frente amplifica a crista a
leste, o que afia o cavado a oeste, que aumenta a advecção de
vorticidade sobre o ciclone\ldots\ realimentação que produz as
\textbf{bombas} (queda $\ge 24$\,hPa/24\,h normalizada por
$\sin\phi/\sin 60°$ — T5), como as ciclogêneses explosivas do
Atlântico Sul ao largo do Uruguai/Argentina.

\section{Síntese da aula}\label{sec:sintese}

\begin{noquadro}[mapa final --- canto do quadro]
\[ \text{baroclinia} \xrightarrow{\text{Eady}} \text{onda que cresce}
\xrightarrow{\text{adv.\ vort.\ + streak}} \text{convergência}
\xrightarrow{-(f+\zeta)D} \text{spin-up} \]
\[ \frac{\partial P_{sfc}}{\partial t} < 0\ \text{(exporta massa)}
\qquad \text{Ekman mata em } \tau_E \qquad
\text{bomba: } 24\,\mathrm{hPa/24h}\times\tfrac{\sin\phi}{\sin60°} \]
\end{noquadro}

Daqui o curso entra na escala convectiva: o que acontece \emph{dentro}
do setor quente instável \veremos{14}{tempestades}.

\section{Exercícios complementares}\label{sec:exercicios}

Soluções (professor) em \texttt{cap13\_solucoes.ipynb}.

\subsection*{Teóricos}

\begin{enumerate}[label=\textbf{T\arabic*.},itemsep=4pt]
  \item A partir das equações do movimento horizontal, deduza o termo
        de esticamento $-(\zeta+f)\nabla\cdot\vec U_h$ da equação da
        vorticidade e interprete-o com conservação de momento angular.
  \item Por análise dimensional com $f$, $N$ e
        $\Lambda = \partial U/\partial z$, monte a taxa de Eady
        $\sigma \sim (f/N)\Lambda$ e a escala
        $L_R = NH/|f|$; avalie $\lambda_{max} = 3{,}9\,L_R$ para a
        troposfera padrão.
  \item Mostre que $\partial P_{sfc}/\partial t =
        -g\int\nabla\cdot(\rho\vec U)\,dz$ e explique por que
        aquecimento diabático da coluna também derruba a pressão à
        superfície.
  \item Derive o vento ageostrófico transversal na entrada/saída de um
        jet streak e monte o padrão de 4 quadrantes para o HS.
  \item Enuncie o critério de Sanders--Gyakum para ``bombas'' com a
        normalização $\sin\phi/\sin 60°$ e calcule o limiar em
        $35°$S e $60°$S.
\end{enumerate}

\subsection*{Numéricos (no notebook)}

\begin{enumerate}[label=\textbf{N\arabic*.},itemsep=4pt]
  \item Varra a amplitude da onda de 500\,hPa e trace $|\zeta|_{max}/|f|$
        e a advecção máxima; onde a hipótese quase-geostrófica começa a
        falhar?
  \item Com o ciclo sazonal de $\partial T/\partial y$, trace o tempo
        de e-dobra de Eady mês a mês e identifique a temporada de
        ciclones do Atlântico Sul.
  \item Desligue a divergência de altitude no dia 3 da EDO de spin-up e
        meça a vida útil do ciclone para $\tau_E = 2, 3, 5$ dias.
  \item Dobre a intensidade do jet streak e meça o máximo de
        divergência e o $w$ estimado por continuidade; compare com o
        levantamento sinótico do Cap.~10.
\end{enumerate}

\vspace{1em}
\noindent\rule{\linewidth}{0.4pt}\\
{\scriptsize Material derivado de R.~Stull, \emph{Practical Meteorology}
(CC BY-NC-SA 4.0); este documento herda a mesma licença.}

\end{document}
